\documentclass[11pt,a4paper]{article}
% preamble.tex — estilo común de todos los apuntes Froth.
% Cada apunte hace % preamble.tex — estilo común de todos los apuntes Froth.
% Cada apunte hace % preamble.tex — estilo común de todos los apuntes Froth.
% Cada apunte hace \input{preamble} justo después de \documentclass.
\usepackage[margin=2.4cm]{geometry}
\usepackage{xcolor}
\definecolor{litblue}{RGB}{31,79,121}
\definecolor{litgreen}{RGB}{27,94,32}
\definecolor{litorange}{RGB}{230,81,0}
\usepackage{parskip}
\usepackage{titlesec}
\titleformat{\section}{\Large\bfseries\color{litblue}}{\thesection}{0.6em}{}
\titleformat{\subsection}{\large\bfseries\color{litblue}}{\thesubsection}{0.6em}{}
\usepackage{tcolorbox}
\usepackage{fancyhdr}
\pagestyle{fancy}
\fancyhf{}
\fancyhead[L]{\small\color{litblue}Froth — Apuntes de ML}
\fancyhead[R]{\small\thepage}
\renewcommand{\headrulewidth}{0.4pt}
\usepackage[colorlinks=true,linkcolor=litblue,urlcolor=litblue]{hyperref}

% Cabecera de cada apunte: \cabecera{numero}{titulo}
\newcommand{\cabecera}[2]{%
  \begin{tcolorbox}[colback=litblue,coltext=white,colframe=litblue,arc=2mm]
    {\Large\bfseries Apunte #1 — #2}\\[3pt]
    {\small Proyecto Froth · aprender Machine Learning construyendo}
  \end{tcolorbox}\vspace{0.8em}}

% Cajas temáticas
\newtcolorbox{concepto}[1]{colback=blue!4,colframe=litblue,title={Concepto: #1},arc=1.5mm}
\newtcolorbox{practica}{colback=green!5,colframe=litgreen,title={Pruébalo tú (teclado en mano)},arc=1.5mm}
\newtcolorbox{ojo}{colback=orange!8,colframe=litorange,title={Ojo / error típico},arc=1.5mm}
\newtcolorbox{resumen}{colback=gray!8,colframe=gray!60,title={En una frase},arc=1.5mm}

\newcommand{\codigo}[1]{\texttt{#1}}
 justo después de \documentclass.
\usepackage[margin=2.4cm]{geometry}
\usepackage{xcolor}
\definecolor{litblue}{RGB}{31,79,121}
\definecolor{litgreen}{RGB}{27,94,32}
\definecolor{litorange}{RGB}{230,81,0}
\usepackage{parskip}
\usepackage{titlesec}
\titleformat{\section}{\Large\bfseries\color{litblue}}{\thesection}{0.6em}{}
\titleformat{\subsection}{\large\bfseries\color{litblue}}{\thesubsection}{0.6em}{}
\usepackage{tcolorbox}
\usepackage{fancyhdr}
\pagestyle{fancy}
\fancyhf{}
\fancyhead[L]{\small\color{litblue}Froth — Apuntes de ML}
\fancyhead[R]{\small\thepage}
\renewcommand{\headrulewidth}{0.4pt}
\usepackage[colorlinks=true,linkcolor=litblue,urlcolor=litblue]{hyperref}

% Cabecera de cada apunte: \cabecera{numero}{titulo}
\newcommand{\cabecera}[2]{%
  \begin{tcolorbox}[colback=litblue,coltext=white,colframe=litblue,arc=2mm]
    {\Large\bfseries Apunte #1 — #2}\\[3pt]
    {\small Proyecto Froth · aprender Machine Learning construyendo}
  \end{tcolorbox}\vspace{0.8em}}

% Cajas temáticas
\newtcolorbox{concepto}[1]{colback=blue!4,colframe=litblue,title={Concepto: #1},arc=1.5mm}
\newtcolorbox{practica}{colback=green!5,colframe=litgreen,title={Pruébalo tú (teclado en mano)},arc=1.5mm}
\newtcolorbox{ojo}{colback=orange!8,colframe=litorange,title={Ojo / error típico},arc=1.5mm}
\newtcolorbox{resumen}{colback=gray!8,colframe=gray!60,title={En una frase},arc=1.5mm}

\newcommand{\codigo}[1]{\texttt{#1}}
 justo después de \documentclass.
\usepackage[margin=2.4cm]{geometry}
\usepackage{xcolor}
\definecolor{litblue}{RGB}{31,79,121}
\definecolor{litgreen}{RGB}{27,94,32}
\definecolor{litorange}{RGB}{230,81,0}
\usepackage{parskip}
\usepackage{titlesec}
\titleformat{\section}{\Large\bfseries\color{litblue}}{\thesection}{0.6em}{}
\titleformat{\subsection}{\large\bfseries\color{litblue}}{\thesubsection}{0.6em}{}
\usepackage{tcolorbox}
\usepackage{fancyhdr}
\pagestyle{fancy}
\fancyhf{}
\fancyhead[L]{\small\color{litblue}Froth — Apuntes de ML}
\fancyhead[R]{\small\thepage}
\renewcommand{\headrulewidth}{0.4pt}
\usepackage[colorlinks=true,linkcolor=litblue,urlcolor=litblue]{hyperref}

% Cabecera de cada apunte: \cabecera{numero}{titulo}
\newcommand{\cabecera}[2]{%
  \begin{tcolorbox}[colback=litblue,coltext=white,colframe=litblue,arc=2mm]
    {\Large\bfseries Apunte #1 — #2}\\[3pt]
    {\small Proyecto Froth · aprender Machine Learning construyendo}
  \end{tcolorbox}\vspace{0.8em}}

% Cajas temáticas
\newtcolorbox{concepto}[1]{colback=blue!4,colframe=litblue,title={Concepto: #1},arc=1.5mm}
\newtcolorbox{practica}{colback=green!5,colframe=litgreen,title={Pruébalo tú (teclado en mano)},arc=1.5mm}
\newtcolorbox{ojo}{colback=orange!8,colframe=litorange,title={Ojo / error típico},arc=1.5mm}
\newtcolorbox{resumen}{colback=gray!8,colframe=gray!60,title={En una frase},arc=1.5mm}

\newcommand{\codigo}[1]{\texttt{#1}}

\begin{document}
\cabecera{32}{h-index por cluster: cuántos papers son ``imperdibles''}

\section{El problema (tu propia orden anti-números-mágicos)}

En la guía de lectura queríamos marcar los must-reads de cada subtópico. ¿Top 3? No:
``top 3'' es un número mágico — igual de arbitrario para un cluster de 110 papers de
geología que para uno de 9 de tomografía. Pediste \emph{investigar cómo investigar}:
un umbral calculado \textbf{por corpus y por cluster}.

\section{La idea robada de la bibliometría: el h-index}

El h-index se inventó para medir investigadores, pero la definición funciona sobre
CUALQUIER lista de conteos de citas:

\begin{concepto}{h-index de un cluster}
El mayor $h$ tal que el cluster tiene $h$ papers con al menos $h$ citas cada uno.
Ordenas por citas descendentes y buscas el punto donde el rango supera a las citas:
justo ahí la evidencia deja de sostenerse a sí misma.
\end{concepto}

Por qué es elegante: \textbf{el propio cluster fija su vara}. Un cluster mega-citado
exige cientos de citas para entrar al must-read; uno joven se conforma con decenas.
Cero parámetros que tunear.

\section{Los datos reales de Beauvoir}

\begin{itemize}
  \item Geología/granitos (110 papers) $\rightarrow$ $h = 35$: literatura madura, la
        vara quedó altísima sola.
  \item Tomografía XCT (9 papers) $\rightarrow$ $h = 3$: nicho joven, vara proporcional.
  \item Informes USGS (mayoría con 0--2 citas) $\rightarrow$ el h-index degenera
        ($h > n/2$, marcaría ``casi todo es imperdible''): ahí entra el \textbf{fallback
        head/tail breaks} — se corta en la media de citas del cluster, que en
        distribuciones de cola larga separa la cabeza pesada del resto ($\rightarrow$ 7).
\end{itemize}

\begin{ojo}
Citas $\neq$ calidad: miden atención acumulada, castigan lo reciente (los 59 papers de
2025--2027 con 0 citas son investigación fresca, no mala). El must-read dice ``esto es
lo canónico que el campo espera que hayas leído'', no ``esto es lo único que vale''.
\end{ojo}

\section{Dónde vive esta regla (y una lección de 2026-08-03)}

Hasta hoy este apunte describía un fallback que \textbf{nunca se había programado}: el
h-index sí estaba en el código, el corte por la media no. Se implementó al construir el
export a Obsidian, en \codigo{froth/export\_obsidian.py} (\codigo{\_must\_reads} y
\codigo{\_mean\_cut}).

\begin{ojo}
\textbf{Un desvío pequeño rompió la regla entera.} Al programarlo, Claude repitió el corte
por la media \emph{hasta seis veces} en vez de hacerlo UNA, creyendo que así sería más
selectivo. Efecto medido sobre tu corpus de tesis: en el subtema de 340 papers hay un
clásico con 4.250 citas que arrastra la media, y al repetir el corte la lista se
estrangulaba hasta \textbf{1 paper}. Con un solo corte da 36. En total, la versión repetida
dejaba con 2 papers o menos a \textbf{20 de los 29 subtemas de 30+ papers}; la regla de
este apunte no lo hace en ninguno.

Lección: cuando un apunte fija un método, el código debe seguirlo \emph{literalmente}.
``Lo hago un poco más estricto'' no es una mejora, es otro método, y hay que medirlo
antes de llamarlo mejor.
\end{ojo}

Números actuales con la regla correcta: corpus de tesis, 3.867 papers en 104 subtemas
$\rightarrow$ 893 must-reads (mediana 6 por subtema, ninguno se queda corto).

\begin{resumen}
Must-reads sin números mágicos: h-index por cluster (auto-calibrado a la madurez de cada
subtópico) con fallback head/tail breaks cuando las citas son demasiado planas para que
el h-index discrimine.
\end{resumen}

\end{document}
